\documentclass[11pt,a4paper]{article}
\usepackage[spanish,es-nodecimaldot]{babel}
\usepackage[utf8]{inputenc}
\usepackage[T1]{fontenc}
\usepackage{lmodern}
\usepackage{amsmath,amssymb,amsthm,mathtools}
\usepackage{geometry}
\usepackage{enumitem}
\usepackage{xcolor}
\usepackage{hyperref}
\usepackage{tikz}
\usepackage{array}
\usepackage{booktabs}
\geometry{margin=2.2cm}
\setlist[itemize]{topsep=3pt,itemsep=2pt}
\setlist[enumerate]{topsep=3pt,itemsep=2pt}
\hypersetup{colorlinks=true, linkcolor=blue!60!black, urlcolor=blue!60!black}

\title{Resumen de Grafos para el Parcial\\\large Temas 1--5}
\author{}
\date{}

\newtheorem*{prop}{Proposición}
\newtheorem*{obs}{Observación}
\newtheorem*{ej}{Ejemplo}
\newtheorem*{dem}{Demostración}

\begin{document}
\maketitle
\tableofcontents
\bigskip

\section{Definiciones y herramientas básicas}

Un \textbf{grafo no dirigido} es un par $G=(V,E)$ donde $V$ es el conjunto de vértices y $E$ el conjunto de aristas. Si las aristas tienen orientación, hablamos de \textbf{dígrafo}. Si tienen peso, de \textbf{grafo ponderado}.

\subsection*{Conceptos básicos}
\begin{itemize}
    \item \textbf{Grado} de un vértice $v$: número de aristas incidentes en $v$. Se denota $\deg(v)$.
    \item \textbf{Recorrido}: sucesión de vértices adyacentes. Puede repetir vértices y aristas.
    \item \textbf{Camino simple}: recorrido sin repetir vértices.
    \item \textbf{Ciclo}: camino cerrado de longitud al menos 3 sin repetir vértices salvo el primero y el último.
    \item \textbf{Conexo}: entre cualquier par de vértices existe un camino.
    \item \textbf{Componente conexa}: subgrafo conexo maximal.
    \item \textbf{Bipartito}: $V$ se puede partir en dos clases $V_1,V_2$ de forma que toda arista une un vértice de $V_1$ con uno de $V_2$.
\end{itemize}

\subsection*{Lema del apretón de manos}
\[
\sum_{v\in V} \deg(v)=2|E|.
\]
En particular, la suma de grados siempre es par.

\begin{dem}
Cada arista contribuye exactamente $2$ a la suma de grados: una vez en cada extremo.
\end{dem}

\subsection*{Consecuencias rápidas}
\begin{itemize}
    \item Un grafo simple con $n$ vértices tiene a lo sumo $\binom{n}{2}$ aristas.
    \item No puede haber un número impar de vértices de grado impar.
\end{itemize}

\subsection*{Isomorfismo: invariantes útiles}
Si dos grafos son isomorfos, entonces tienen:
\begin{itemize}
    \item el mismo número de vértices y aristas,
    \item la misma multiconjunto de grados,
    \item la misma cantidad de ciclos de cada longitud,
    \item la misma cantidad de recorridos de longitud fija,
    \item la misma conectividad y la misma bipartición, si aplica.
\end{itemize}
Estos criterios no siempre prueban isomorfismo, pero sí permiten descartar muchos casos.

\section{Matrices del grafo}

\subsection{Matriz de adyacencia}
Para un grafo simple no dirigido de $n$ vértices,
\[
A=(a_{ij}),\qquad a_{ij}=\begin{cases}
1,&\text{si } i\sim j,\\
0,&\text{si no.}
\end{cases}
\]
En grafos simples, $A$ es simétrica y su diagonal es nula.

\subsection{Matriz de incidencia}
Si el grafo tiene vértices $v_1,\dots,v_n$ y aristas $e_1,\dots,e_m$, la matriz de incidencia es
\[
M=(m_{ij}),\qquad m_{ij}=\begin{cases}
1,&\text{si } v_i \text{ incide en } e_j,\\
0,&\text{si no.}
\end{cases}
\]
En un grafo simple no dirigido, cada columna tiene exactamente dos unos.

\subsection{Relación entre ambas}
Si $D=\operatorname{diag}(\deg(v_1),\dots,\deg(v_n))$, entonces
\[
MM^T=A+D.
\]

\begin{dem}
La entrada $(i,j)$ de $MM^T$ es el producto escalar de la fila $i$ con la fila $j$ de $M$.
\begin{itemize}
    \item Si $i\neq j$, cuenta cuántas aristas unen $v_i$ con $v_j$: eso es $a_{ij}$.
    \item Si $i=j$, cuenta cuántas aristas inciden en $v_i$: eso es $\deg(v_i)$.
\end{itemize}
Luego $MM^T$ coincide con $A$ fuera de la diagonal y con $D$ en la diagonal.
\end{dem}

\subsection{Potencias de la matriz de adyacencia}
La entrada $(i,j)$ de $A^k$ da el número de \textbf{recorridos} de longitud $k$ entre $i$ y $j$.

\begin{ej}
Si
\[
A=\begin{pmatrix}0&1\\1&0\end{pmatrix},\qquad
A^2=\begin{pmatrix}1&0\\0&1\end{pmatrix},
\]
entonces hay un recorrido de longitud $2$ que empieza y acaba en cada vértice: ir y volver por la única arista.
\end{ej}

Consecuencias importantes:
\begin{itemize}
    \item $\operatorname{tr}(A^2)=2|E|$.
    \item $\operatorname{tr}(A^3)=6\cdot$(número de triángulos).
    \item Si $A+A^2+\cdots +A^{n-1}$ no tiene entradas nulas, el grafo es conexo.
\end{itemize}

\section{Árboles}

\subsection{Caracterizaciones}
Un \textbf{árbol} es un grafo conexo sin ciclos. Equivalentemente:
\begin{itemize}
    \item entre cada pareja de vértices hay un único camino simple,
    \item es conexo y tiene $|E|=|V|-1$,
    \item es acíclico y tiene $|E|=|V|-1$.
\end{itemize}

\subsection{Propiedad fundamental}
Si $T$ es un árbol con $n$ vértices, entonces
\[
|E(T)|=n-1.
\]

\begin{dem}
Por inducción en $n$. Para $n=1$ es trivial. Si $n\ge 2$, un árbol tiene una hoja. Al eliminarla se obtiene un árbol con $n-1$ vértices y una arista menos. Por hipótesis inductiva, el árbol reducido tiene $(n-1)-1=n-2$ aristas, así que el original tiene $n-1$.
\end{dem}

\subsection{Todo árbol con $n\ge 2$ tiene al menos dos hojas}

\begin{dem}
Tomemos un camino simple de longitud máxima
\[
v_0,v_1,\dots,v_k.
\]
Si $v_0$ tuviera otro vecino distinto de $v_1$, o bien ese vecino ya está en el camino y se formaría un ciclo, o bien no está y podríamos alargar el camino. Ambas cosas contradicen la definición de árbol o la maximalidad del camino. Luego $\deg(v_0)=1$. Igual para $v_k$.
\end{dem}

\subsection{Puentes en árboles}
Toda arista de un árbol es un puente: si eliminas una arista, el árbol deja de ser conexo.

\subsection{Árboles enraizados}
Un árbol enraizado es un árbol con un vértice distinguido como \textbf{raíz}. Para listarlos con pocos vértices:
\begin{enumerate}
    \item enumera primero los árboles no isomorfos,
    \item luego coloca la raíz en vértices no equivalentes por simetría.
\end{enumerate}

\begin{ej}
Con $5$ vértices solo hay dos árboles no isomorfos: la estrella $K_{1,4}$ y el árbol con sucesión de grados $(1,1,1,2,3)$. Enraizando en vértices no equivalentes se obtienen $6$ árboles enraizados no isomorfos.
\end{ej}

\subsection{Código de Prüfer}
Para un árbol etiquetado con vértices $\{1,\dots,n\}$, el \textbf{código de Prüfer} es una secuencia de longitud $n-2$ construida así:
\begin{enumerate}
    \item toma la hoja de menor etiqueta,
    \item escribe la etiqueta de su vecino,
    \item elimina la hoja,
    \item repite hasta que queden dos vértices.
\end{enumerate}

\begin{ej}
Para el árbol con aristas
\[
1-3,\;2-3,\;3-4,\;4-5,\;4-6,
\]
el código es
\[
(3,3,4,4).
\]
\end{ej}

Propiedad útil:
\[
\text{un vértice } v \text{ aparece } \deg(v)-1 \text{ veces en el código de Prüfer.}
\]

\section{Planaridad y fórmula de Euler}

Un grafo es \textbf{plano} si puede dibujarse en el plano sin cruces de aristas, y \textbf{planar} si admite algún dibujo plano.

\subsection{Fórmula de Euler}
Si $G$ es plano y conexo, entonces
\[
|V|-|E|+|F|=2,
\]
donde $F$ es el número de caras, incluyendo la exterior.

\begin{ej}
Si un grafo plano conexo tiene $11$ aristas y $8$ caras, entonces
\[
|V|=2+|E|-|F|=2+11-8=5.
\]
Si además fuera simple, sería imposible, porque un grafo simple con $5$ vértices tiene como máximo $10$ aristas.
\end{ej}

\subsection{Cota planar simple}
Si $G$ es planar simple y tiene $n\ge 3$ vértices, entonces
\[
|E|\le 3n-6.
\]
Si además no tiene triángulos, entonces
\[
|E|\le 2n-4.
\]

\begin{dem}
En un grafo plano simple cada cara tiene longitud al menos $3$, así que
\[
3|F|\le 2|E|,
\]
porque cada arista cuenta dos veces al recorrer las caras. Sustituyendo en Euler:
\[
2=|V|-|E|+|F|\le |V|-|E|+\frac{2|E|}{3}=|V|-\frac{|E|}{3}.
\]
Reordenando,
\[
|E|\le 3|V|-6.
\]
\end{dem}

\section{Recorridos del grafo: DFS y BFS}

\subsection{DFS (Depth First Search)}
Explora ``en profundidad'': desde un vértice, va al vecino no visitado de menor etiqueta (si se fija ese criterio), y retrocede cuando no puede seguir.

Usos principales:
\begin{itemize}
    \item construir un árbol recubridor,
    \item comprobar conectividad,
    \item contar componentes conexas,
    \item detectar ciclos,
    \item buscar puentes y vértices de corte.
\end{itemize}

\subsection{BFS (Breadth First Search)}
Explora ``por niveles'': primero todos los vecinos del origen, luego los vecinos de éstos, etc.

Usos principales:
\begin{itemize}
    \item distancias mínimas en grafos no ponderados,
    \item excentricidad, radio y diámetro,
    \item caminos más cortos en número de aristas.
\end{itemize}

\subsection{Comparación}
\begin{center}
\begin{tabular}{>{\bfseries}p{3cm}p{5.3cm}p{5.3cm}}
\toprule
 & DFS & BFS\\
\midrule
Idea & Baja por una rama y retrocede & Explora por capas\\
Estructura & Muy bueno para puentes, ciclos, componentes & Muy bueno para distancias\\
Distancias mínimas & No en general & Sí, si no hay pesos\\
Estructura de datos & Pila / recursión & Cola\\
\bottomrule
\end{tabular}
\end{center}

\subsection{Puentes con DFS}
Una arista es un \textbf{puente} si al eliminarla aumenta el número de componentes conexas.

Con DFS se usan tiempos de entrada $tin[v]$ y valores $low[v]$:
\begin{itemize}
    \item $tin[v]$: instante en que visitamos $v$ por primera vez,
    \item $low[v]$: menor tiempo de entrada accesible desde el subárbol de $v$ usando aristas del árbol DFS y, como mucho, una arista de retroceso.
\end{itemize}

La arista del árbol DFS $(v,to)$ es puente si y solo si
\[
low[to]>tin[v].
\]

\begin{dem}
Si $low[to]\le tin[v]$, desde el subárbol de $to$ se puede volver a $v$ o a un antepasado suyo por otro camino; por tanto, al quitar $(v,to)$ el grafo sigue conectado en esa zona. Si $low[to]>tin[v]$, no hay vuelta hacia arriba y el subárbol de $to$ queda colgando únicamente de esa arista.
\end{dem}

\section{Distancias, excentricidad, radio y diámetro}

La \textbf{distancia} $d(u,v)$ es la longitud del camino más corto entre $u$ y $v$.

La \textbf{excentricidad} de $v$ es
\[
e(v)=\max_{u\in V} d(v,u).
\]
El \textbf{radio} es
\[
\operatorname{rad}(G)=\min_{v\in V} e(v),
\]
y el \textbf{diámetro}
\[
\operatorname{diam}(G)=\max_{v\in V} e(v).
\]

En grafos no ponderados se calculan con BFS:
\begin{itemize}
    \item distancia entre dos vértices: BFS desde uno de ellos,
    \item excentricidad de $v$: BFS desde $v$ y tomar la mayor distancia,
    \item radio y diámetro: BFS desde todos los vértices.
\end{itemize}

\section{Árboles recubridores mínimos}

En un grafo ponderado, un \textbf{árbol recubridor mínimo} (MST) es un árbol recubridor cuya suma de pesos es mínima.

\subsection{Kruskal}
Idea: ordenar aristas por peso creciente e ir añadiendo la menor que no forme ciclo.

\textbf{Esquema:}
\begin{enumerate}
    \item Ordena las aristas por peso.
    \item Recorre la lista de menor a mayor.
    \item Añade una arista si conecta dos componentes distintas.
    \item Para cuando tengas $n-1$ aristas.
\end{enumerate}

\subsection{Prim}
Idea: empezar en un vértice y crecer siempre por la arista más barata que conecta el árbol actual con un vértice nuevo.

\textbf{Esquema:}
\begin{enumerate}
    \item Empieza en cualquier vértice.
    \item Entre las aristas que salen del árbol actual, escoge la de menor peso que no forme ciclo.
    \item Repite hasta tener todos los vértices.
\end{enumerate}

\subsection{Cuándo usar cada uno}
\begin{itemize}
    \item \textbf{Kruskal}: muy natural cuando te dan una tabla o lista de aristas.
    \item \textbf{Prim}: muy cómodo cuando el grafo ya está dibujado y quieres crecer desde una zona.
\end{itemize}

\section{Dijkstra}

Sirve para hallar distancias mínimas desde un origen en grafos con pesos \textbf{no negativos}.

\subsection*{Idea}
Mantiene un conjunto de vértices cuya distancia definitiva ya se conoce. En cada paso fija el vértice no procesado con menor distancia provisional y relaja sus vecinos.

\textbf{Esquema:}
\begin{enumerate}
    \item Inicializa $dist(s)=0$ y $dist(v)=\infty$ para $v\neq s$.
    \item Elige el vértice no fijado de distancia mínima.
    \item Actualiza a sus vecinos con la regla
    \[
    dist(v)\leftarrow \min\{dist(v),dist(u)+w(u,v)\}.
    \]
    \item Repite.
\end{enumerate}

\begin{obs}
Dijkstra no sirve con pesos negativos.
\end{obs}

\section{Tema 3: ALT}

\subsection{A* y heurísticas}
A* selecciona el vértice con menor valor
\[
f(n)=g(n)+h(n),
\]
donde $g(n)$ es el coste ya recorrido y $h(n)$ es una estimación del coste restante.

Para garantizar optimalidad, la heurística debe ser \textbf{admisible}:
\[
h(n)\le d(n,\text{goal}).
\]

\subsection{Heurística ALT}
ALT significa \textit{A*, Landmarks, Triangle inequality}. Se eligen landmarks $\ell_1,\dots,\ell_k$ y se precomputan distancias. La heurística es
\[
h(n)=\max_{\ell\in L} \bigl|dist(\ell,goal)-dist(\ell,n)\bigr|.
\]

\textbf{Por qué es admisible:} por desigualdad triangular,
\[
|d(\ell,g)-d(\ell,n)|\le d(n,g).
\]
Tomar el máximo de cotas inferiores sigue dando una cota inferior.

\section{Ford--Fulkerson, flujo máximo y corte mínimo}

\subsection{Conceptos}
En una red dirigida con fuente $s$ y sumidero $t$:
\begin{itemize}
    \item cada arista $(u,v)$ tiene una capacidad $c(u,v)$,
    \item un flujo $f(u,v)$ cumple $0\le f(u,v)\le c(u,v)$,
    \item en los vértices intermedios se conserva el flujo.
\end{itemize}

\subsection{Ford--Fulkerson}
Busca \textbf{caminos aumentantes} en la red residual y aumenta el flujo tanto como permita el cuello de botella del camino.

\textbf{Esquema:}
\begin{enumerate}
    \item Empieza con flujo $0$.
    \item Construye la red residual.
    \item Busca un camino $s\leadsto t$.
    \item Aumenta el flujo en el mínimo residual del camino.
    \item Repite hasta que no haya camino aumentante.
\end{enumerate}

\subsection{Teorema max-flow = min-cut}
El valor del flujo máximo es igual a la capacidad del corte mínimo.

\textbf{Cómo sacar el corte mínimo al final del algoritmo:}
\begin{enumerate}
    \item en la red residual final, toma $S$ = vértices alcanzables desde $s$,
    \item toma $T=V\setminus S$,
    \item las aristas del grafo original que van de $S$ a $T$ forman un corte mínimo.
\end{enumerate}

\textbf{Idea intuitiva:} si ya no hay camino aumentante, las aristas que separan la parte alcanzable de la no alcanzable están saturadas y forman el cuello de botella definitivo.

\subsection{Grafos no dirigidos}
También puede aplicarse a grafos no dirigidos interpretando cada arista como un canal por el que el flujo puede circular en ambos sentidos, o modelando adecuadamente su versión dirigida con red residual.

\section{Teoría espectral de grafos}

\subsection{Espectro de un grafo}
El \textbf{espectro} de $G$ es el conjunto de autovalores de su matriz de adyacencia $A$, contando multiplicidades.

\subsection{Espectros conocidos}
\begin{itemize}
    \item Para el grafo completo $K_n$:
    \[
    \operatorname{Spec}(K_n)=\{n-1,-1,\dots,-1\}.
    \]
    \item Para el ciclo $C_n$:
    \[
    \lambda_k=2\cos\left(\frac{2\pi k}{n}\right),\qquad k=0,1,\dots,n-1.
    \]
\end{itemize}

\subsection{Bipartición y simetría del espectro}
Un grafo es bipartito si y solo si su espectro es \textbf{simétrico respecto a $0$}. Esto significa:
\[
\lambda\in \operatorname{Spec}(G) \Longrightarrow -\lambda\in \operatorname{Spec}(G)
\]
con la misma multiplicidad.

\textbf{Ojo:} esto \emph{no} significa que $0$ tenga que ser autovalor. Solo significa que los autovalores aparecen en parejas opuestas.

\begin{ej}
$C_6$ es bipartito y su espectro es simétrico: $\{2,1,1,-1,-1,-2\}$.
\end{ej}

\subsection{Laplaciana}
La matriz laplaciana es
\[
L=D-A,
\]
donde $D$ es la matriz diagonal de grados.

Propiedades básicas:
\begin{itemize}
    \item $L$ es simétrica.
    \item $0$ siempre es autovalor de $L$ con autovector $\mathbf{1}$.
    \item La multiplicidad de $0$ coincide con el número de componentes conexas.
\end{itemize}

\subsection{Autovalor y autovector de Fiedler}
Se llama \textbf{valor de Fiedler} a $\lambda_2$, el segundo menor autovalor de la laplaciana. Su autovector asociado se llama \textbf{vector de Fiedler}.

Interpretación:
\begin{itemize}
    \item si $\lambda_2$ es pequeño, el grafo tiene un cuello de botella y está cerca de desconectarse,
    \item el signo del vector de Fiedler sugiere una partición natural en comunidades,
    \item valores cercanos a $0$ en el autovector suelen corresponder a nodos frontera o puente.
\end{itemize}

\subsection{Conductancia}
La conductancia de un subconjunto $S$ es, intuitivamente, la cantidad de aristas que salen de $S$ relativa al volumen de $S$. Una formulación típica es
\[
h(S)=\frac{|\partial S|}{\min\{\operatorname{vol}(S),\operatorname{vol}(\overline{S})\}}.
\]
La conductancia del grafo es el mínimo sobre cortes no triviales.

\subsection{Cota de Cheeger}
Para la laplaciana normalizada,
\[
\frac{h(G)^2}{2}\le \mu_2\le 2h(G).
\]
Interpretación práctica:
\begin{itemize}
    \item $\mu_2$ pequeño $\Rightarrow$ existe un corte bueno que casi separa el grafo,
    \item $\mu_2$ grande $\Rightarrow$ el grafo está bien mezclado y no tiene cuellos de botella claros.
\end{itemize}

\subsection{Centralidad de intermediación (betweenness)}
La \textbf{betweenness centrality} de un vértice mide cuántos caminos mínimos entre otros pares de vértices pasan por él. Intuitivamente detecta nodos ``intermediarios'' o puentes funcionales.

\section{Coloración de grafos}

\subsection{Número cromático}
El \textbf{número cromático} $\chi(G)$ es el mínimo número de colores necesarios para colorear los vértices de modo que vértices adyacentes tengan colores distintos.

\subsection{Casos clásicos}
\begin{itemize}
    \item Grafo completo:
    \[
    \chi(K_n)=n.
    \]
    \item Ciclo:
    \[
    \chi(C_n)=\begin{cases}
    2,& n \text{ par},\\
    3,& n \text{ impar}.
    \end{cases}
    \]
    \item Rueda $W_n$ (ciclo $C_n$ más centro unido a todos):
    \[
    \chi(W_n)=\begin{cases}
    3,& n \text{ par},\\
    4,& n \text{ impar}.
    \end{cases}
    \]
    \item Árbol con al menos una arista: $\chi(T)=2$.
\end{itemize}

\begin{dem}[Por qué $\chi(C_n)$ vale $2$ o $3$]
Si $n$ es par, se colorea alternando dos colores. Si $n$ es impar, al alternar dos colores el último vértice queda adyacente al primero con el mismo color, así que hacen falta al menos $3$. Con tres colores sí se puede.
\end{dem}

\subsection{Teorema de Brooks}
Sea $\Delta$ el grado máximo de un grafo conexo $G$. Entonces:
\[
\chi(G)\le \Delta,
\]
salvo cuando $G$ es un grafo completo o un ciclo impar, en cuyo caso
\[
\chi(G)=\Delta+1.
\]

\textbf{Idea de uso en examen:} no suele dar el número cromático exacto, pero sí muy buenas cotas.

\subsection{Algoritmos de coloración}
\begin{itemize}
    \item \textbf{Voraz}: recorre vértices en cierto orden y asigna a cada uno el menor color disponible.
    \item \textbf{Welsh--Powell}: ordena por grado decreciente y luego aplica voraz.
    \item \textbf{Brélaz / DSATUR}: en cada paso elige el vértice con mayor saturación, es decir, el que ``ve'' más colores distintos entre sus vecinos ya coloreados.
\end{itemize}

\textbf{Idea práctica:} DSATUR suele dar muy buenas coloraciones y es exacto en ciertos casos como bipartitos y ciclos.

\subsection{Índice cromático}
El \textbf{índice cromático} es el mínimo número de colores para colorear las \textbf{aristas}, de modo que aristas adyacentes tengan colores distintos.

\section{Polinomio cromático}

El polinomio cromático $P_G(k)$ cuenta el número de coloraciones propias de $G$ usando $k$ colores.

\subsection{Ejemplos básicos}
\begin{itemize}
    \item Grafo vacío con $n$ vértices:
    \[
    P_G(k)=k^n.
    \]
    \item Grafo completo $K_n$:
    \[
    P_{K_n}(k)=k(k-1)(k-2)\cdots (k-n+1).
    \]
    \item Árbol con $n$ vértices:
    \[
    P_T(k)=k(k-1)^{n-1}.
    \]
\end{itemize}

\begin{dem}[Polinomio cromático de un árbol]
Elegimos el color de la raíz: $k$ posibilidades. Cada nuevo vértice solo debe evitar el color de su padre, así que tiene $k-1$ opciones. Como un árbol con $n$ vértices tiene $n-1$ aristas,
\[
P_T(k)=k(k-1)^{n-1}.
\]
\end{dem}

\subsection{Propiedades importantes}
Si $P_G(k)$ es de grado $n$, entonces:
\begin{itemize}
    \item el grado es $n=|V|$,
    \item el término independiente es $0$,
    \item los signos de los coeficientes alternan,
    \item el coeficiente de $k^{n-1}$ cambiado de signo es $|E|$,
    \item el grado del último término no nulo es el número de componentes conexas.
\end{itemize}

\subsection{Deleción--contracción}
Si $e$ es una arista de $G$, entonces
\[
P_G(k)=P_{G-e}(k)-P_{G/e}(k).
\]
Aquí:
\begin{itemize}
    \item $G-e$ es el grafo al eliminar la arista $e$,
    \item $G/e$ es el grafo al contraer $e$, identificando sus extremos.
\end{itemize}

\begin{dem}
En $G-e$, los extremos de $e$ o bien reciben colores distintos, y entonces la coloración es válida en $G$, o bien reciben el mismo color, y entonces corresponde a una coloración de $G/e$. Por tanto,
\[
P_{G-e}(k)=P_G(k)+P_{G/e}(k).
\]
Reordenando,
\[
P_G(k)=P_{G-e}(k)-P_{G/e}(k).
\]
\end{dem}

\section{Plantilla mental para resolver ejercicios}

\subsection*{Existencia de grafos}
\begin{itemize}
    \item Usa suma de grados $=2|E|$.
    \item Si es simple, recuerda $|E|\le \binom{n}{2}$.
    \item Si es planar, prueba Euler y las cotas $|E|\le 3n-6$ o $2n-4$.
\end{itemize}

\subsection*{Árboles}
\begin{itemize}
    \item Piensa siempre en $|E|=|V|-1$.
    \item Si el enunciado dice ``ruta única entre cada pareja'', es un árbol.
    \item Un árbol con $n\ge 2$ tiene al menos dos hojas.
\end{itemize}

\subsection*{DFS / BFS}
\begin{itemize}
    \item DFS para estructura: ciclos, puentes, componentes.
    \item BFS para distancias, radio, diámetro, excentricidad.
\end{itemize}

\subsection*{MST}
\begin{itemize}
    \item Kruskal: ordena aristas.
    \item Prim: crece desde un árbol parcial.
\end{itemize}

\subsection*{Coloración}
\begin{itemize}
    \item Busca cliques para cotas inferiores.
    \item Usa bipartición si no hay ciclos impares.
    \item Usa Brooks para cotas superiores.
\end{itemize}

\subsection*{Espectral}
\begin{itemize}
    \item En bipartitos, el espectro es simétrico respecto a $0$.
    \item El segundo autovalor laplaciano detecta conectividad fina y comunidades.
\end{itemize}

\section{Errores típicos}
\begin{itemize}
    \item Confundir ``camino'' con ``recorrido'' al interpretar $A^k$.
    \item Pensar que DFS siempre da distancias mínimas: no, eso lo hace BFS.
    \item Aplicar Dijkstra con pesos negativos.
    \item Creer que ``espectro simétrico respecto a 0'' significa que $0$ es autovalor: no necesariamente.
    \item En polinomio cromático, olvidar que se cuentan coloraciones con vértices etiquetados.
    \item En planaridad, usar Euler sin comprobar conectividad o sin contar la cara exterior.
\end{itemize}

\section{Mini lista de resultados para memorizar}
\[
\sum \deg(v)=2|E|,
\qquad
|E(T)|=|V(T)|-1,
\qquad
|V|-|E|+|F|=2.
\]
\[
MM^T=A+D,
\qquad
(A^k)_{ij}=\#\{\text{recorridos de longitud }k\text{ de }i\text{ a }j\}.
\]
\[
\chi(C_n)=\begin{cases}2,&n\text{ par}\\3,&n\text{ impar}\end{cases}
\qquad
\chi(W_n)=\begin{cases}3,&n\text{ par}\\4,&n\text{ impar}\end{cases}
\]
\[
P_T(k)=k(k-1)^{n-1},
\qquad
P_{K_n}(k)=k(k-1)\cdots(k-n+1).
\]
\[
\text{Puente }(v,to) \iff low[to]>tin[v].
\]
\[
\text{bipartito }\iff \text{espectro simétrico respecto a }0.
\]

\bigskip
\begin{center}
\Large\textbf{Suerte en el parcial.}
\end{center}

\end{document}
